% Options for packages loaded elsewhere
\PassOptionsToPackage{unicode}{hyperref}
\PassOptionsToPackage{hyphens}{url}
\PassOptionsToPackage{dvipsnames,svgnames,x11names}{xcolor}
%
\documentclass[
  11pt,
  a4paper,
]{ltjsarticle}
\usepackage{amsmath,amssymb}
\usepackage{iftex}
\ifPDFTeX
  \usepackage[T1]{fontenc}
  \usepackage[utf8]{inputenc}
  \usepackage{textcomp} % provide euro and other symbols
\else % if luatex or xetex
  \usepackage{unicode-math} % this also loads fontspec
  \defaultfontfeatures{Scale=MatchLowercase}
  \defaultfontfeatures[\rmfamily]{Ligatures=TeX,Scale=1}
\fi
\usepackage{lmodern}
\ifPDFTeX\else
  % xetex/luatex font selection
\fi
% Use upquote if available, for straight quotes in verbatim environments
\IfFileExists{upquote.sty}{\usepackage{upquote}}{}
\IfFileExists{microtype.sty}{% use microtype if available
  \usepackage[]{microtype}
  \UseMicrotypeSet[protrusion]{basicmath} % disable protrusion for tt fonts
}{}
\makeatletter
\@ifundefined{KOMAClassName}{% if non-KOMA class
  \IfFileExists{parskip.sty}{%
    \usepackage{parskip}
  }{% else
    \setlength{\parindent}{0pt}
    \setlength{\parskip}{6pt plus 2pt minus 1pt}}
}{% if KOMA class
  \KOMAoptions{parskip=half}}
\makeatother
\usepackage{xcolor}
\usepackage[margin=24mm]{geometry}
\setlength{\emergencystretch}{3em} % prevent overfull lines
\providecommand{\tightlist}{%
  \setlength{\itemsep}{0pt}\setlength{\parskip}{0pt}}
\setcounter{secnumdepth}{5}
\ifLuaTeX
  \usepackage{selnolig}  % disable illegal ligatures
\fi
\IfFileExists{bookmark.sty}{\usepackage{bookmark}}{\usepackage{hyperref}}
\IfFileExists{xurl.sty}{\usepackage{xurl}}{} % add URL line breaks if available
\urlstyle{same}
\hypersetup{
  colorlinks=true,
  linkcolor={blue},
  filecolor={Maroon},
  citecolor={Blue},
  urlcolor={blue},
  pdfcreator={LaTeX via pandoc}}

\author{}
\date{}

\begin{document}

\hypertarget{gauss-ux6d88ux53bbux6cd5}{%
\section{Gauss 消去法}\label{gauss-ux6d88ux53bbux6cd5}}

Gauss
消去法は、基本行変形によって連立方程式を解きやすい形に変える方法です。

許される操作は

\begin{enumerate}
\def\labelenumi{\arabic{enumi}.}
\tightlist
\item
  行の交換
\item
  非零定数倍
\item
  他の行の定数倍を加える
\end{enumerate}

です。

\hypertarget{ux306aux305cux89e3ux96c6ux5408ux304cux5909ux308fux3089ux306aux3044ux306eux304b}{%
\subsection{なぜ解集合が変わらないのか}\label{ux306aux305cux89e3ux96c6ux5408ux304cux5909ux308fux3089ux306aux3044ux306eux304b}}

これらはすべて「同値な方程式へ置き換える操作」です。行列だけを機械的に変形しているのではなく、同じ制約集合を別の形で記述しています。

\hypertarget{ux672cux8cea}{%
\subsection{本質}\label{ux672cux8cea}}

消去法の目的は数値を出すことだけではありません。pivot と free variable
を露出させることで

\begin{itemize}
\tightlist
\item
  ランク
\item
  解の一意性
\item
  零空間の次元
\end{itemize}

を同時に明らかにします。

\end{document}
